\makeatletter
\def\input@path{{./}{../}{../../}{../../../}{../../../../}{preamble/}{../preamble/}{../../preamble/}{../../../preamble/}}
\makeatother
\input{preamble_exams.tex}
\title{Grade 11 -- Term 1 \\ C9 Classroom Activity}
\author{}\date{}
\pagestyle{empty}
\linespread{1.04}
\setlength{\parskip}{3pt}
\begin{document}
	\newgeometry{left=0.4in,right=0.4in,top=0.35in,bottom=0.35in}
	\begin{center}{\large\bfseries G11 T1 C9 -- Student's t versus Normal}\end{center}
	\textbf{Models and simplifying approximation (all pages).} Each question supplies a continuous model for a signed deviation. Compare $X_T=m+aT$, $T\sim t_\nu$, with $X_N=m+aZ$, $Z\sim N(0,1)$, using the same location $m$ and scale $a>0$. Here $\boxed{\nu=n-1}$, where $n$ is the number of observations in the sample. For simplicity, use the stated ``approximate standard deviation'' directly as $a$. This is a teaching approximation: $\operatorname{Var}(X_T)=a^2\frac{\nu}{\nu-2}$ and $\operatorname{SD}(X_T)=a\sqrt{\frac{\nu}{\nu-2}}$ for $\nu>2$, whereas $\operatorname{Var}(X_N)=a^2$. The approximation is coarser for small $\nu$ and improves as $\nu$ increases. Use the supplied models without adjusting $a$. Negative deviations mean below the reference; positive deviations mean above it.

	\textbf{Reference-table method.} Use the supplied normal CDF table and Student's $t$ critical-value table. In row $\nu$ and column $\alpha$, $P(T>t^*)=\alpha$; by symmetry, $P(T<-t^*)=\alpha$. Standardize using $c=(x-m)/a$ for either model; back-transform using $x=m+ac$. Read $t^*$ to three decimals. Direct bounds retain enough digits to recover the intended $t$ entries. For normal direct probabilities, round $c$ to two decimals before lookup. For normal inverse problems, choose the closest CDF entry.

	\textbf{Comparison.} For each requested probability or cutoff, calculate both results, then
	\[
	E_A=|\text{$t$-result}-\text{$N$-result}|,\qquad
	E_{Pt}=\frac{100E_A}{|\text{$t$-result}|},\qquad
	E_{Pn}=\frac{100E_A}{|\text{$N$-result}|}.
	\]
	Report $E_{Pt}$ and $E_{Pn}$ as percentages; a zero denominator makes the corresponding percentage undefined. For inverse intervals, compare the lower and upper cutoffs separately. Report probabilities and their $E_A$ to four decimals, contextual cutoffs and their $E_A$ to two decimals, and percentage errors to two decimals. All results are table-based approximations.

	\textbf{Six worked examples}
	\begin{enumerate}[leftmargin=*,label=\textbf{\arabic*.},itemsep=5pt]
		\item \textbf{Direct right tail -- cost error.} The cost deviation from budget of 12 batches, in dollars, has mean $-10$ and approximate standard deviation $20$. Use $X_T=-10+20T$, $T\sim t_{11}$, and $X_N=-10+20Z$. Find $P(X>17.26)$ under each model, then $E_A$, $E_{Pt}$, and $E_{Pn}$.
		\par Standardize: $c=\frac{17.26+10}{20}=1.363$. Row $11$, $\alpha=0.10$: $p_T=0.1000$. For normal, $z\approx1.36$, so $p_N=1-\Phi(1.36)=1-0.9131=0.0869$.
		$E_A\approx0.0131$, $E_{Pt}\approx13.10\%$, $E_{Pn}\approx15.07\%$. About $10\%$ under $t$, compared with $8.69\%$ under normal, exceed budget by more than $\$17.26$.
		\item \textbf{Direct left tail -- packing error.} Let $X$ be a package's mass deviation from target, in grams. A sample of $11$ packages has mean $m=1$ and approximate standard deviation $a=4$, so $\nu=10$. Use $X_T=1+4T$, $T\sim t_{10}$, and $X_N=1+4Z$, $Z\sim N(0,1)$. What is the probability that a package is more than $4.488$ grams below target? Find both probabilities and all three errors.
		\par \textit{Region and standardization:} below target means a negative deviation. We need the area to the left of $-4.488$, namely $P(X<-4.488)$. Standardize: $c=\frac{-4.488-1}{4}=-1.372$.
		\par \textit{Cumulative probabilities:} write $F_{10}(c)=P(T<c)$ and $\Phi(c)=P(Z<c)$. Row $10$, $\alpha=0.10$ gives $P(T>1.372)=0.10$; symmetry gives $p_T=F_{10}(-1.372)\approx0.1000$. For normal, round $c$ to $-1.37$: $p_N=\Phi(-1.37)=1-\Phi(1.37)=1-0.9147=0.0853$.
		\par \textit{Comparison:} $E_A=|0.1000-0.0853|=0.0147$, $E_{Pt}=100(0.0147)/0.1000=14.70\%$, $E_{Pn}=100(0.0147)/0.0853\approx17.23\%$. About $10\%$ under $t$, compared with $8.53\%$ under normal, have a mass more than $4.49$ grams below target.
		\item \textbf{Direct central interval -- delivery error.} Let $X$ be a delivery's arrival-time deviation from schedule, in minutes. A sample of $11$ deliveries has mean $m=2$ and approximate standard deviation $a=5$, so $\nu=10$. Use $X_T=2+5T$, $T\sim t_{10}$, and $X_N=2+5Z$, $Z\sim N(0,1)$. Find the probability of arrival between $4.86$ minutes early and $11.06$ minutes late under each model, then all three errors.
		\par \textit{Region and standardization:} we need the area between the bounds, $P(-4.86<X<11.06)$. This interval crosses the mean but is not symmetric about it. Standardize both endpoints:
		\[
		c_L=\frac{-4.86-2}{5}=-1.372,\qquad c_U=\frac{11.06-2}{5}=1.812.
		\]
		\par \textit{Cumulative subtraction:} row $10$ gives $P(T>1.372)=0.10$ and $P(T>1.812)=0.05$. Thus $F_{10}(-1.372)\approx0.10$ by symmetry and $F_{10}(1.812)\approx1-0.05=0.95$. Subtract the cumulative area at the lower bound from that at the upper bound:
		\[
		p_T=F_{10}(1.812)-F_{10}(-1.372)\approx0.9500-0.1000=0.8500.
		\]
		For normal, round to $z_L=-1.37$, $z_U=1.81$. The table gives $\Phi(-1.37)=1-0.9147=0.0853$ and $\Phi(1.81)=0.9649$, so
		\[
		p_N=\Phi(1.81)-\Phi(-1.37)=0.9649-0.0853=0.8796.
		\]
		\par \textit{Comparison:} $E_A=|0.8500-0.8796|=0.0296$, $E_{Pt}=100(0.0296)/0.8500\approx3.48\%$, $E_{Pn}=100(0.0296)/0.8796\approx3.37\%$. About $85\%$ under $t$, compared with $87.96\%$ under normal, arrive within these limits.
		
		
%		\item \textbf{Inverse right tail -- electricity error.} Let $X$ be a day's electricity-use deviation from forecast, in kWh. A sample of $11$ days has mean $m=2$ and approximate standard deviation $a=5$, so $\nu=10$. Use $X_T=2+5T$, $T\sim t_{10}$, and $X_N=2+5Z$, $Z\sim N(0,1)$. Find the deviation cutoff exceeded on only $5\%$ of days under each model, then all three errors.
%		\par \textit{Region and cumulative target:} the required right tail is $P(X>x)=0.05$. Convert to the left cumulative probability: $P(X<x)=1-0.05=0.95$. Standardize the unknown cutoff: $c=(x-2)/5$, so seek $F_{10}(c_T)=0.95$ and $\Phi(c_N)=0.95$.
%		\par \textit{Table lookup:} the $t$ table uses right tails, so $\alpha=1-0.95=0.05$; row $10$ gives $c_T=1.812$. The normal table entries $\Phi(1.64)=0.9495$ and $\Phi(1.65)=0.9505$ are equally close to $0.95$; use $c_N=1.64$ by the stated tie rule.
%		\par \textit{Back-transformation:} $x=m+ac$ gives $x_T=2+5(1.812)=\boxed{11.06\ \text{kWh}}$ and $x_N=2+5(1.64)=\boxed{10.20\ \text{kWh}}$.
%		\par \textit{Comparison:} $E_A=|11.06-10.20|=0.86$ kWh, $E_{Pt}=100(0.86)/11.06\approx7.78\%$, $E_{Pn}=100(0.86)/10.20\approx8.43\%$. Approximately $5\%$ of days exceed forecast by more than $11.06$ kWh under $t$, or $10.20$ kWh under normal.
		
%		\item \textbf{Inverse left tail -- sales error.} Let $X$ be a week's sales deviation from target, in thousands of dollars. A sample of $11$ weeks has mean $m=-1$ and approximate standard deviation $a=3$, so $\nu=10$. Use $X_T=-1+3T$, $T\sim t_{10}$, and $X_N=-1+3Z$, $Z\sim N(0,1)$. Find the deviation cutoff below which $2.5\%$ of weeks fall under each model, then all three errors.
%		\par \textit{Region and cumulative target:} the required left tail is $P(X<x)=0.025$. Since this is already a left-tail probability, the cumulative target remains $0.025$; its complement is $P(X>x)=0.975$. Standardize: $c=(x+1)/3$, so seek $F_{10}(c_T)=0.025$ and $\Phi(c_N)=0.025$.
%		\par \textit{Table lookup:} row $10$, $\alpha=0.025$ gives $P(T>2.228)=0.025$. By symmetry, $P(T<-2.228)=0.025$, so $c_T=-2.228$. For normal, use $\Phi(1.96)=0.9750$ and $\Phi(-1.96)=1-0.9750=0.0250$, giving $c_N=-1.96$.
%		\par \textit{Back-transformation:} $x=-1+3c$ gives $x_T=-1+3(-2.228)=-7.684\approx\boxed{-7.68\ \text{thousand dollars}}$ and $x_N=-1+3(-1.96)=\boxed{-6.88\ \text{thousand dollars}}$.
%		\par \textit{Comparison:} use the unrounded cutoffs: $E_A=|-7.684-(-6.88)|=0.804\approx0.80$ thousand dollars, $E_{Pt}=100(0.804)/7.684\approx10.46\%$, $E_{Pn}=100(0.804)/6.88\approx11.69\%$. About $2.5\%$ of weeks have sales more than $7.68$ thousand dollars below target under $t$, or $6.88$ thousand dollars below target under normal.
%		\item \textbf{Inverse central interval -- demand error.} The demand deviation from forecast over 14 weeks, in tonnes, has mean $5$ and approximate standard deviation $10$. Use $X_T=5+10T$, $T\sim t_{13}$, and $X_N=5+10Z$. Find the central $90\%$ limits under each model, then all three errors for each endpoint.
%		\par Standardize: $c_L=(L-5)/10$, $c_U=(U-5)/10$. Each tail is $0.05$: row $13$ gives $(c_{T,L},c_{T,U})=(-1.771,1.771)$; normal gives $(-1.64,1.64)$. Back-transform: $(L_T,U_T)=(5-17.71,5+17.71)=(-12.71,22.71)$ tonnes; $(L_N,U_N)=(5-16.40,5+16.40)=(-11.40,21.40)$ tonnes.
%		\par Lower: $E_A\approx1.31$, $E_{Pt}\approx10.31\%$, $E_{Pn}\approx11.49\%$. Upper: $E_A\approx1.31$, $E_{Pt}\approx5.77\%$, $E_{Pn}\approx6.12\%$. The normal interval is narrower.

	\end{enumerate}
	\textbf{PROBLEMS}\\
	For each problem, show the table lookup, standardization and, where needed, back-transformation. Find the same requested probability or cutoff under both models, then calculate $E_A$, $E_{Pt}$, and $E_{Pn}$ and interpret the results in context. For inverse intervals, calculate the three errors for each endpoint separately. Work on separate paper.
	\begin{enumerate}[leftmargin=*,label=\textbf{\arabic*.},itemsep=5pt,parsep=0pt,topsep=2pt]
		\item \textbf{Bus arrivals.} The arrival-time deviation from schedule of 7 buses, in minutes, has mean $2$ and approximate standard deviation $6$. Compare $X_T=2+6T$, $T\sim t_{6}$, with $X_N=2+6Z$. Find the probability that the deviation is below $-9.658$, i.e., $P(X<-9.658)$, and interpret the threshold. Use both models, then calculate $E_A$, $E_{Pt}$, and $E_{Pn}$.
		\item \textbf{Shipping costs.} The cost deviation of 16 shipments, in thousand USD, has mean $-15$ and approximate standard deviation $40$. Compare $X_T=-15+40T$, $T\sim t_{15}$, with $X_N=-15+40Z$. Find the probability that the deviation exceeds $-49.64$, i.e., $P(X>-49.64)$, and explain why it exceeds $0.5$. Use both models, then calculate $E_A$, $E_{Pt}$, and $E_{Pn}$.
		\item \textbf{Water demand.} The water-use deviation from forecast over 20 days, in cubic metres, has mean $10$ and approximate standard deviation $20$. Compare $X_T=10+20T$, $T\sim t_{19}$, with $X_N=10+20Z$. Find the probability that the deviation lies between $27.22$ and $44.58$, i.e., $P(27.22<X<44.58)$. Use both models, then calculate $E_A$, $E_{Pt}$, and $E_{Pn}$.
		\item \textbf{Harvest forecasts.} The harvested-mass deviation from forecast of 9 harvests, in tonnes, has mean $-2$ and approximate standard deviation $4$. Compare $X_T=-2+4T$, $T\sim t_{8}$, with $X_N=-2+4Z$. Find the probability that the deviation lies between $-11.224$ and $-5.556$, i.e., $P(-11.224<X<-5.556)$. Use both models, then calculate $E_A$, $E_{Pt}$, and $E_{Pn}$.
		\item \textbf{Weekly revenue alerts.} The revenue deviation from plan over 5 weeks, in thousands of dollars, has mean $1$ and approximate standard deviation $3$. Compare $X_T=1+3T$, $T\sim t_{4}$, with $X_N=1+3Z$. An alert occurs when the deviation is below $-5.396$ or above $5.599$. Find the alert probability, i.e., $P(X<-5.396\text{ or }X>5.599)$. Use both models, then calculate $E_A$, $E_{Pt}$, and $E_{Pn}$.
		\item \textbf{Cold-storage sensors.} The reading deviation from reference of 11 sensors, in degrees Celsius, has mean $-2$ and approximate standard deviation $6$. Compare $X_T=-2+6T$, $T\sim t_{10}$, with $X_N=-2+6Z$. Find the cutoff $x$ below which $10\%$ of deviations fall, i.e., $P(X<x)=0.10$. Use both models, then calculate $E_A$, $E_{Pt}$, and $E_{Pn}$.
		\item \textbf{Rail journey delays.} The journey-time deviation from schedule of 24 rail journeys, in minutes, has mean $3$ and approximate standard deviation $8$. Compare $X_T=3+8T$, $T\sim t_{23}$, with $X_N=3+8Z$. Find the cutoff $x$ exceeded by only $2.5\%$ of journeys, i.e., $P(X>x)=0.025$. Use both models, then calculate $E_A$, $E_{Pt}$, and $E_{Pn}$.
		\item \textbf{Electricity forecasting.} The electricity-use deviation from forecast over 19 weeks, in kWh, has mean $10$ and approximate standard deviation $40$. Compare $X_T=10+40T$, $T\sim t_{18}$, with $X_N=10+40Z$. Find limits $L<U$ enclosing the central $90\%$ of deviations, i.e., $P(L<X<U)=0.90$, with equal excluded tails. Use both models, then calculate $E_A$, $E_{Pt}$, and $E_{Pn}$ (for each cutoff separately).
		\item \textbf{Packaging inspection.} The mass deviation from target of 6 packages, in grams, has mean $-0.5$ and approximate standard deviation $1.5$. Compare $X_T=-0.5+1.5T$, $T\sim t_{5}$, with $X_N=-0.5+1.5Z$. Find limits $L<U$ leaving $5\%$ of deviations outside, equally split between the tails, i.e., $P(X<L)=P(X>U)=0.025$. Use both models, then calculate $E_A$, $E_{Pt}$, and $E_{Pn}$ (for each cutoff separately).
		\item \textbf{Processing-time alerts.} The processing-time deviation from target of 13 batches, in minutes, has mean $1$ and approximate standard deviation $5$. Compare $X_T=1+5T$, $T\sim t_{12}$, with $X_N=1+5Z$. Find limits $L<U$ leaving $15\%$ of deviations outside, with $5\%$ below $L$ and $10\%$ above $U$, i.e., $P(X<L)=0.05$ and $P(X>U)=0.10$. Use both models, then calculate $E_A$, $E_{Pt}$, and $E_{Pn}$ (for each cutoff separately).
		
		\newpage
		
		\item \textbf{Hospital waiting.} The waiting-time deviation from schedule of 13 appointments, in minutes, has mean $1$ and approximate standard deviation $5$. Compare $X_T=1+5T$, $T\sim t_{12}$, with $X_N=1+5Z$. Find the cutoff $x$ below which $10\%$ of deviations fall, i.e., $P(X<x)=0.10$. Use both models, then calculate $E_A$, $E_{Pt}$, and $E_{Pn}$.
		
		\item \textbf{Production tolerance.} The production-output deviation from target of 21 runs, in units, has mean $0$ and approximate standard deviation $12$. Compare $X_T=0+12T$, $T\sim t_{20}$, with $X_N=0+12Z$. Find limits $L<U$ enclosing the central $95\%$ of deviations, i.e., $P(L<X<U)=0.95$, with equal excluded tails. Use both models, then calculate $E_A$, $E_{Pt}$, and $E_{Pn}$ (for each cutoff separately).
		\item \textbf{Electricity peak.} The peak-demand deviation from forecast over 16 days, in MW, has mean $3$ and approximate standard deviation $7$. Compare $X_T=3+7T$, $T\sim t_{15}$, with $X_N=3+7Z$. Find the cutoff $x$ exceeded by $1\%$ of deviations, i.e., $P(X>x)=0.01$. Use both models, then calculate $E_A$, $E_{Pt}$, and $E_{Pn}$.
		\item \textbf{Transit reliability.} The arrival-time deviation from schedule of 9 journeys, in minutes, has mean $-1$ and approximate standard deviation $3$. Compare $X_T=-1+3T$, $T\sim t_{8}$, with $X_N=-1+3Z$. Find alert limits $L<U$ with $5\%$ below $L$ and $2.5\%$ above $U$, i.e., $P(X<L)=0.05$ and $P(X>U)=0.025$. Use both models, then calculate $E_A$, $E_{Pt}$, and $E_{Pn}$ (for each cutoff separately).
		\item \textbf{Shipping weight.} The shipment-mass deviation from target of 11 shipments, in kg, has mean $2$ and approximate standard deviation $4$. Compare $X_T=2+4T$, $T\sim t_{10}$, with $X_N=2+4Z$. Find limits $L<U$ enclosing the central $98\%$ of deviations, i.e., $P(L<X<U)=0.98$, with equal excluded tails. Use both models, then calculate $E_A$, $E_{Pt}$, and $E_{Pn}$ (for each cutoff separately).
	\end{enumerate}
	
	\textbf{\large Solutions}\\
Use the supplied approximation $X_T=m+aT$, $T\sim t_\nu$, and $X_N=m+aZ$, with $\nu=n-1$: the stated $a$ is used directly as the scale. Write $F_\nu(c)=P(T<c)$ and $\Phi(c)=P(Z<c)$. Since both models are continuous, including or excluding an endpoint does not change a probability.

Read positive $t$ critical values to three decimals from the row for $\nu$ and right-tail column $\alpha$: the entry $t^*$ means $P(T>t^*)=\alpha$. These printed critical values approximate the underlying distribution. In the calculations below, the associated table probabilities are used as the working values. Equalities in subsequent arithmetic are exact for those selected table values; the resulting model probabilities and cutoffs remain approximations.

For direct normal calculations, first round each standardized boundary to two decimals, then read the CDF to four decimals. For inverse normal calculations, choose the closest printed CDF entry. At cumulative targets $0.0500$ and $0.9500$, the printed entries at magnitudes $1.64$ and $1.65$ are equally close: consistently choose $|z|=1.64$. Thus $\Phi(1.64)=0.9495$ is used for a $0.9500$ target, and $\Phi(-1.64)=1-0.9495=0.0505$ for a $0.0500$ target.

Comparison errors measure differences between the two model results, not errors against observed truth. Use displayed table probabilities for direct errors and unrounded back-transformed cutoffs for inverse errors. Keep absolute values in percentage denominators; a zero denominator would make that percentage undefined. Report probabilities and their absolute differences to four decimals, contextual cutoffs and their absolute differences to two decimals with units, and percentage errors to two decimals.
\begingroup
\begin{enumerate}[leftmargin=*,label=\textbf{\arabic*.},itemsep=12pt,parsep=3pt,topsep=5pt]
\item \begin{minipage}[t]{\linewidth}\textbf{Bus arrivals ($\nu=6$).}\nopagebreak
\par\textbf{Step 1: Information and event.}\nopagebreak\enspace $n=7$, $m=2$, $a=6$, so $\nu=n-1=7-1=6$. The quantity is arrival-time deviation in minutes. A negative deviation means early arrival; the event is a left tail.\[
P(X<-9.658).
\]

\end{minipage}\par
\par\noindent\begin{minipage}{\linewidth}
\par\textbf{Step 2: Standardize.}\nopagebreak\enspace \[
c=\frac{x-m}{a}.
\]
Subtract the model location and divide by its positive scale; this preserves the direction of the inequalities.
\[
c=\frac{-9.658-2}{6}=\frac{-11.658}{6}=-1.943.
\]

\end{minipage}\par
\par\noindent\begin{minipage}{\linewidth}
\par\textbf{Step 3: Student's $t$ probability.}\nopagebreak\enspace In row $6$, column $\alpha=0.05$, the entry is $t^*=1.943$:
\[
P(T>1.943)\approx0.0500.
\]
Symmetry gives the same area to the left of the negative critical value:
\[
p_T=P(T<-1.943)=P(T>1.943)\approx\boxed{0.0500}.
\]

\end{minipage}\par
\par\textbf{Step 4: Normal probability.}\nopagebreak\enspace Round the standardized boundary: $-1.943\to z=-1.94$. The table gives $\Phi(1.94)=0.9738$. Use symmetry to find the negative cumulative area:
\[
\Phi(-z)=1-\Phi(z).
\]
\[
p_N=\Phi(-1.94)=1-\Phi(1.94)=1-0.9738=\boxed{0.0262}.
\]

\par\textbf{Step 5: Comparison errors.}\nopagebreak\enspace \par\noindent\begin{minipage}{\linewidth}
\par\textit{Bus arrivals comparison.}
\[
E_A=|r_T-r_N|,\qquad E_{Pt}=\frac{100E_A}{|r_T|},\qquad E_{Pn}=\frac{100E_A}{|r_N|}.
\]
Here $r_T$ and $r_N$ are the displayed probabilities being compared.
\[
\begin{aligned}
E_A&=|0.0500-(0.0262)|=0.0238,\\
E_{Pt}&=\frac{100|0.0500-(0.0262)|}{|0.0500|}=\frac{100(0.0238)}{0.0500}\approx 47.60\%,\\
E_{Pn}&=\frac{100|0.0500-(0.0262)|}{|0.0262|}=\frac{100(0.0238)}{0.0262}\approx 90.84\%.
\end{aligned}
\]
\end{minipage}\par

\par\noindent\begin{minipage}{\linewidth}
\par\textbf{Step 6: Interpretation.}\nopagebreak\enspace The threshold means arriving more than $9.658$ minutes early (about $9.66$ minutes). The $t$ model assigns about $5.00\%$ to this event, and the normal model about $2.62\%$. The normal probability is lower by $0.0238$, or $2.38$ percentage points.
\end{minipage}\par
\item \begin{minipage}[t]{\linewidth}\textbf{Shipping costs ($\nu=15$).}\nopagebreak
\par\textbf{Step 1: Information and event.}\nopagebreak\enspace $n=16$, $m=-15$, $a=40$, so $\nu=n-1=16-1=15$. Cost deviation is in thousand USD. The negative threshold is $49.64$ thousand USD below budget. We need the area to its right, including less negative and positive deviations.\[
P(X>-49.64).
\]

\end{minipage}\par
\par\noindent\begin{minipage}{\linewidth}
\par\textbf{Step 2: Standardize.}\nopagebreak\enspace \[
c=\frac{x-m}{a}.
\]
Subtract the model location and divide by its positive scale; this preserves the direction of the inequalities.
\[
c=\frac{-49.64-(-15)}{40}=\frac{-49.64+15}{40}=\frac{-34.64}{40}=-0.866.
\]

\end{minipage}\par
\par\noindent\begin{minipage}{\linewidth}
\par\textbf{Step 3: Student's $t$ probability.}\nopagebreak\enspace Row $15$, column $\alpha=0.20$, gives $t^*=0.866$:
\[
P(T>0.866)\approx0.2000.
\]
Symmetry transfers this to the left tail; taking its complement gives the requested right tail:
\[
\begin{aligned}P(T<-0.866)&=P(T>0.866)\approx0.2000,\\p_T=P(T>-0.866)&=1-P(T<-0.866)\\&\approx1-0.2000=\boxed{0.8000}.\end{aligned}
\]

\end{minipage}\par
\par\textbf{Step 4: Normal probability.}\nopagebreak\enspace Round $-0.866\to z=-0.87$. The table gives $\Phi(0.87)=0.8078$. First use symmetry, then the right-tail complement:
\[
\Phi(-z)=1-\Phi(z),\qquad P(Z>z)=1-\Phi(z).
\]
\[
\begin{aligned}\Phi(-0.87)&=1-0.8078=0.1922,\\p_N=P(Z>-0.87)&=1-\Phi(-0.87)\\&=1-0.1922=\boxed{0.8078}.\end{aligned}
\]

\par\textbf{Step 5: Comparison errors.}\nopagebreak\enspace \par\noindent\begin{minipage}{\linewidth}
\par\textit{Shipping costs comparison.}
\[
E_A=|r_T-r_N|,\qquad E_{Pt}=\frac{100E_A}{|r_T|},\qquad E_{Pn}=\frac{100E_A}{|r_N|}.
\]
Here $r_T$ and $r_N$ are the displayed probabilities being compared.
\[
\begin{aligned}
E_A&=|0.8000-(0.8078)|=0.0078,\\
E_{Pt}&=\frac{100|0.8000-(0.8078)|}{|0.8000|}=\frac{100(0.0078)}{0.8000}\approx 0.98\%,\\
E_{Pn}&=\frac{100|0.8000-(0.8078)|}{|0.8078|}=\frac{100(0.0078)}{0.8078}\approx 0.97\%.
\end{aligned}
\]
\end{minipage}\par

\par\noindent\begin{minipage}{\linewidth}
\par\textbf{Step 6: Interpretation.}\nopagebreak\enspace The chance of a cost deviation exceeding $-49.64$ thousand USD is about $80.00\%$ under $t$ and $80.78\%$ under normal. Both exceed $0.5$ because $-49.64<-15=m$: each model is symmetric about $-15$, so the area to the right of a threshold below that centre includes the right half plus some of the left half. The normal probability is higher by $0.0078$, or $0.78$ percentage points.
\end{minipage}\par
\item \begin{minipage}[t]{\linewidth}\textbf{Water demand ($\nu=19$).}\nopagebreak
\par\textbf{Step 1: Information and event.}\nopagebreak\enspace $n=20$, $m=10$, $a=20$, so $\nu=n-1=20-1=19$. Water-use deviation is in cubic metres. Both positive bounds mean use above forecast, and both bounds are above the model location $10$.\[
P(27.22<X<44.58).
\]

\end{minipage}\par
\par\noindent\begin{minipage}{\linewidth}
\par\textbf{Step 2: Standardize.}\nopagebreak\enspace \[
c=\frac{x-m}{a}.
\]
Subtract the model location and divide by its positive scale; this preserves the direction of the inequalities.
\[
\begin{aligned}c_L&=\frac{27.22-10}{20}=\frac{17.22}{20}=0.861,\\c_U&=\frac{44.58-10}{20}=\frac{34.58}{20}=1.729.\end{aligned}
\]

\end{minipage}\par
\par\noindent\begin{minipage}{\linewidth}
\par\textbf{Step 3: Student's $t$ probability.}\nopagebreak\enspace Row $19$ gives $t^*=0.861$ in column $\alpha=0.20$, and $t^*=1.729$ in column $\alpha=0.05$. Thus
\[
P(T>0.861)\approx0.2000,\qquad P(T>1.729)\approx0.0500.
\]
Each left cumulative area is one minus its right tail. Subtracting the lower cumulative area removes values below the interval:
\[
P(c_L<T<c_U)=F_\nu(c_U)-F_\nu(c_L).
\]
\[
\begin{aligned}F_{19}(0.861)&\approx1-0.2000=0.8000,\\F_{19}(1.729)&\approx1-0.0500=0.9500,\\p_T&\approx0.9500-0.8000=\boxed{0.1500}.\end{aligned}
\]

\end{minipage}\par
\par\textbf{Step 4: Normal probability.}\nopagebreak\enspace Round each boundary separately: $0.861\to z_L=0.86$ and $1.729\to z_U=1.73$. The CDF entries are $\Phi(0.86)=0.8051$ and $\Phi(1.73)=0.9582$. Subtract cumulative areas:
\[
P(z_L<Z<z_U)=\Phi(z_U)-\Phi(z_L).
\]
\[
p_N=\Phi(1.73)-\Phi(0.86)=0.9582-0.8051=\boxed{0.1531}.
\]

\par\textbf{Step 5: Comparison errors.}\nopagebreak\enspace \par\noindent\begin{minipage}{\linewidth}
\par\textit{Water demand comparison.}
\[
E_A=|r_T-r_N|,\qquad E_{Pt}=\frac{100E_A}{|r_T|},\qquad E_{Pn}=\frac{100E_A}{|r_N|}.
\]
Here $r_T$ and $r_N$ are the displayed probabilities being compared.
\[
\begin{aligned}
E_A&=|0.1500-(0.1531)|=0.0031,\\
E_{Pt}&=\frac{100|0.1500-(0.1531)|}{|0.1500|}=\frac{100(0.0031)}{0.1500}\approx 2.07\%,\\
E_{Pn}&=\frac{100|0.1500-(0.1531)|}{|0.1531|}=\frac{100(0.0031)}{0.1531}\approx 2.02\%.
\end{aligned}
\]
\end{minipage}\par

\par\noindent\begin{minipage}{\linewidth}
\par\textbf{Step 6: Interpretation.}\nopagebreak\enspace Water use is between $27.22$ and $44.58$ cubic metres above forecast with probability about $15.00\%$ under $t$ and $15.31\%$ under normal. The normal probability is higher by $0.0031$, or $0.31$ percentage points.
\end{minipage}\par
\item \begin{minipage}[t]{\linewidth}\textbf{Harvest forecasts ($\nu=8$).}\nopagebreak
\par\textbf{Step 1: Information and event.}\nopagebreak\enspace $n=9$, $m=-2$, $a=4$, so $\nu=n-1=9-1=8$. Harvested-mass deviation is in tonnes. Both negative bounds represent shortfalls below forecast, and both lie below the model location $-2$.\[
P(-11.224<X<-5.556).
\]

\end{minipage}\par
\par\noindent\begin{minipage}{\linewidth}
\par\textbf{Step 2: Standardize.}\nopagebreak\enspace \[
c=\frac{x-m}{a}.
\]
Subtract the model location and divide by its positive scale; this preserves the direction of the inequalities.
\[
\begin{aligned}c_L&=\frac{-11.224-(-2)}{4}=\frac{-11.224+2}{4}=\frac{-9.224}{4}=-2.306,\\c_U&=\frac{-5.556-(-2)}{4}=\frac{-5.556+2}{4}=\frac{-3.556}{4}=-0.889.\end{aligned}
\]

\end{minipage}\par
\par\noindent\begin{minipage}{\linewidth}
\par\textbf{Step 3: Student's $t$ probability.}\nopagebreak\enspace In row $8$, column $\alpha=0.025$ gives $t^*=2.306$, and column $\alpha=0.20$ gives $t^*=0.889$:
\[
P(T>2.306)\approx0.0250,\qquad P(T>0.889)\approx0.2000.
\]
By symmetry these are the cumulative areas at the corresponding negative bounds. Use cumulative subtraction:
\[
P(c_L<T<c_U)=F_\nu(c_U)-F_\nu(c_L).
\]
\[
\begin{aligned}F_8(-2.306)&=P(T>2.306)\approx0.0250,\\F_8(-0.889)&=P(T>0.889)\approx0.2000,\\p_T&\approx0.2000-0.0250=\boxed{0.1750}.\end{aligned}
\]

\end{minipage}\par
\par\textbf{Step 4: Normal probability.}\nopagebreak\enspace Round $-2.306\to z_L=-2.31$ and $-0.889\to z_U=-0.89$. The positive CDF entries are $\Phi(2.31)=0.9896$ and $\Phi(0.89)=0.8133$. Use symmetry before subtracting:
\[
\Phi(-z)=1-\Phi(z),\qquad P(z_L<Z<z_U)=\Phi(z_U)-\Phi(z_L).
\]
\[
\begin{aligned}\Phi(-2.31)&=1-0.9896=0.0104,\\\Phi(-0.89)&=1-0.8133=0.1867,\\p_N&=0.1867-0.0104=\boxed{0.1763}.\end{aligned}
\]

\par\textbf{Step 5: Comparison errors.}\nopagebreak\enspace \par\noindent\begin{minipage}{\linewidth}
\par\textit{Harvest forecasts comparison.}
\[
E_A=|r_T-r_N|,\qquad E_{Pt}=\frac{100E_A}{|r_T|},\qquad E_{Pn}=\frac{100E_A}{|r_N|}.
\]
Here $r_T$ and $r_N$ are the displayed probabilities being compared.
\[
\begin{aligned}
E_A&=|0.1750-(0.1763)|=0.0013,\\
E_{Pt}&=\frac{100|0.1750-(0.1763)|}{|0.1750|}=\frac{100(0.0013)}{0.1750}\approx 0.74\%,\\
E_{Pn}&=\frac{100|0.1750-(0.1763)|}{|0.1763|}=\frac{100(0.0013)}{0.1763}\approx 0.74\%.
\end{aligned}
\]
\end{minipage}\par

\par\noindent\begin{minipage}{\linewidth}
\par\textbf{Step 6: Interpretation.}\nopagebreak\enspace A shortfall between $5.556$ and $11.224$ tonnes (about $5.56$ to $11.22$ tonnes below forecast) has probability about $17.50\%$ under $t$ and $17.63\%$ under normal. The normal probability is higher by $0.0013$, or $0.13$ percentage points.
\end{minipage}\par
\item \begin{minipage}[t]{\linewidth}\textbf{Weekly revenue alerts ($\nu=4$).}\nopagebreak
\par\textbf{Step 1: Information and event.}\nopagebreak\enspace $n=5$, $m=1$, $a=3$, so $\nu=n-1=5-1=4$. Revenue deviation is in thousands of dollars. The negative limit marks revenue below plan; the positive limit marks revenue above plan. An alert is in either of two disjoint tails.\[
P(X<-5.396\text{ or }X>5.599).
\]

\end{minipage}\par
\par\noindent\begin{minipage}{\linewidth}
\par\textbf{Step 2: Standardize.}\nopagebreak\enspace \[
c=\frac{x-m}{a}.
\]
Subtract the model location and divide by its positive scale; this preserves the direction of the inequalities.
\[
\begin{aligned}c_L&=\frac{-5.396-1}{3}=\frac{-6.396}{3}=-2.132,\\c_U&=\frac{5.599-1}{3}=\frac{4.599}{3}=1.533.\end{aligned}
\]

\end{minipage}\par
\par\noindent\begin{minipage}{\linewidth}
\par\textbf{Step 3: Student's $t$ probability.}\nopagebreak\enspace In row $4$, column $\alpha=0.05$ gives $t^*=2.132$, and column $\alpha=0.10$ gives $t^*=1.533$:
\[
P(T>2.132)\approx0.0500,\qquad P(T>1.533)\approx0.1000.
\]
Symmetry gives $P(T<-2.132)=P(T>2.132)$. A value cannot be in both alert tails, so their probabilities are added:
\[
P(T<c_L\text{ or }T>c_U)=F_\nu(c_L)+[1-F_\nu(c_U)].
\]
\[
p_T\approx0.0500+0.1000=\boxed{0.1500}.
\]

\end{minipage}\par
\par\textbf{Step 4: Normal probability.}\nopagebreak\enspace Round $-2.132\to z_L=-2.13$ and $1.533\to z_U=1.53$. The table gives $\Phi(2.13)=0.9834$ and $\Phi(1.53)=0.9370$. Use symmetry for the left tail and a complement for the right:
\[
\Phi(-z)=1-\Phi(z),\qquad P(Z>z)=1-\Phi(z).
\]
\[
\begin{aligned}\Phi(-2.13)&=1-0.9834=0.0166,\\P(Z>1.53)&=1-0.9370=0.0630.\end{aligned}
\]
Add these disjoint tail areas:
\[
\begin{aligned}p_N&=\Phi(z_L)+[1-\Phi(z_U)]\\&=0.0166+0.0630=\boxed{0.0796}.\end{aligned}
\]

\par\textbf{Step 5: Comparison errors.}\nopagebreak\enspace \par\noindent\begin{minipage}{\linewidth}
\par\textit{Weekly revenue alerts comparison.}
\[
E_A=|r_T-r_N|,\qquad E_{Pt}=\frac{100E_A}{|r_T|},\qquad E_{Pn}=\frac{100E_A}{|r_N|}.
\]
Here $r_T$ and $r_N$ are the displayed probabilities being compared.
\[
\begin{aligned}
E_A&=|0.1500-(0.0796)|=0.0704,\\
E_{Pt}&=\frac{100|0.1500-(0.0796)|}{|0.1500|}=\frac{100(0.0704)}{0.1500}\approx 46.93\%,\\
E_{Pn}&=\frac{100|0.1500-(0.0796)|}{|0.0796|}=\frac{100(0.0704)}{0.0796}\approx 88.44\%.
\end{aligned}
\]
\end{minipage}\par

\par\noindent\begin{minipage}{\linewidth}
\par\textbf{Step 6: Interpretation.}\nopagebreak\enspace An alert means revenue more than $5.396$ thousand dollars below plan or more than $5.599$ thousand dollars above plan (about $5.40$ below or $5.60$ above). Its probability is about $15.00\%$ under $t$ and $7.96\%$ under normal. The normal probability is lower by $0.0704$, or $7.04$ percentage points.
\end{minipage}\par
\item \begin{minipage}[t]{\linewidth}\textbf{Cold-storage sensors ($\nu=10$).}\nopagebreak
\par\textbf{Step 1: Information and required cutoffs.}\nopagebreak\enspace $n=11$, $m=-2$, $a=6$, so $\nu=n-1=11-1=10$. The quantity is a sensor reading deviation in degrees Celsius; negative values mean below the reference reading.\[
P(X<x)=0.10.
\]
This is already a left-tail target: $P(X<x)=0.1000$; its complement is $P(X>x)=1-0.10=0.90$.
\end{minipage}\par
\par\noindent\begin{minipage}{\linewidth}
\par\textbf{Step 2: Standardize the unknowns.}\nopagebreak\enspace \[
c=\frac{x-m}{a}.
\]
\[
c=\frac{x-(-2)}{6}=\frac{x+2}{6}.
\]
A negative standardized cutoff lies below the model location $m$; a positive one lies above it. The numerical standardized values are found from the probability targets before back-transformation.
\end{minipage}\par
\par\noindent\begin{minipage}{\linewidth}
\par\textbf{Step 3: Student's $t$ cutoffs.}\nopagebreak\enspace \par\textit{Cutoff.} Seek $F_{10}(c_{T})=0.1000$. This is below $0.5$, so the cutoff is negative. By symmetry use right-tail column $\alpha=0.1$. Row $10$, column $\alpha=0.1$, gives $t^*=1.372$:\[
P(T>1.372)\approx0.1000.
\]
\[
F_{10}(-1.372)=P(T>1.372)\approx0.1000,\qquad c_{T}=-1.372.
\]

\end{minipage}\par
\par\textbf{Step 4: Normal cutoffs.}\nopagebreak\enspace \par\noindent\begin{minipage}{\linewidth}
\par\textit{Cutoff.} The left cumulative target is $\Phi(c_{N})=0.1000$. It requires a negative $z$ because the target is below $0.5$. For a negative lookup use $\Phi(-z)=1-\Phi(z)$.\[
\Phi(-1.28)=1-\Phi(1.28)=1-0.8997=0.1003.
\]
The adjacent entry $\Phi(-1.29)=1-0.9015=0.0985$ is farther from $0.1000$.\[
c_{N}=-1.28.
\]
\end{minipage}\par

\par\noindent\begin{minipage}{\linewidth}
\par\textbf{Step 5: Back-transform to contextual units.}\nopagebreak\enspace \[
x=m+ac.
\]
Multiply each selected standardized cutoff by the scale, then add the location. Retain the unrounded result for the comparison errors.
\par\textit{Cutoff back-transformation.}
\[
\begin{aligned}x_T&=-2+6(-1.372)=-2+(-8.232)\\&=-10.232\approx\boxed{-10.23\ \text{degrees Celsius}},\\x_N&=-2+6(-1.28)=-2+(-7.68)\\&=-9.68\approx\boxed{-9.68\ \text{degrees Celsius}}.\end{aligned}
\]

\end{minipage}\par
\par\textbf{Step 6: Comparison errors.}\nopagebreak\enspace \par\noindent\begin{minipage}{\linewidth}
\par\textit{Cutoff comparison.}
\[
E_A=|r_T-r_N|,\qquad E_{Pt}=\frac{100E_A}{|r_T|},\qquad E_{Pn}=\frac{100E_A}{|r_N|}.
\]
Here $r_T$ and $r_N$ are the unrounded contextual cutoffs being compared.
\[
\begin{aligned}
E_A&=|-10.232-(-9.68)|=0.552\approx 0.55\ \text{degrees Celsius},\\
E_{Pt}&=\frac{100|-10.232-(-9.68)|}{|-10.232|}=\frac{100(0.552)}{10.232}\approx 5.39\%,\\
E_{Pn}&=\frac{100|-10.232-(-9.68)|}{|-9.68|}=\frac{100(0.552)}{9.68}\approx 5.70\%.
\end{aligned}
\]
\end{minipage}\par

\par\noindent\begin{minipage}{\linewidth}
\par\textbf{Step 7: Interpretation.}\nopagebreak\enspace About $10\%$ of reading deviations lie below $-10.23$ degrees Celsius under $t$, or below $-9.68$ degrees Celsius under normal. These are deviations below the reference, not absolute storage temperatures. The normal cutoff is closer to the mean deviation by about $0.55$ degrees Celsius.
\end{minipage}\par
\item \begin{minipage}[t]{\linewidth}\textbf{Rail journey delays ($\nu=23$).}\nopagebreak
\par\textbf{Step 1: Information and required cutoffs.}\nopagebreak\enspace $n=24$, $m=3$, $a=8$, so $\nu=n-1=24-1=23$. The quantity is journey-time deviation in minutes; positive values mean journeys longer than scheduled.\[
P(X>x)=0.025.
\]
Convert the right tail to a left cumulative target before lookup:\[
P(X<x)=1-P(X>x)=1-0.025=0.9750.
\]

\end{minipage}\par
\par\noindent\begin{minipage}{\linewidth}
\par\textbf{Step 2: Standardize the unknowns.}\nopagebreak\enspace \[
c=\frac{x-m}{a}.
\]
\[
c=\frac{x-(3)}{8}.
\]
A negative standardized cutoff lies below the model location $m$; a positive one lies above it. The numerical standardized values are found from the probability targets before back-transformation.
\end{minipage}\par
\par\noindent\begin{minipage}{\linewidth}
\par\textbf{Step 3: Student's $t$ cutoffs.}\nopagebreak\enspace \par\textit{Cutoff.} Seek $F_{23}(c_{T})=0.9750$. This is above $0.5$, so the cutoff is positive. The right-tail column is $\alpha=1-0.9750=0.025$. Row $23$, column $\alpha=0.025$, gives $t^*=2.069$:\[
P(T>2.069)\approx0.0250.
\]
\[
F_{23}(2.069)=1-P(T>2.069)\approx1-0.0250=0.9750.
\]
\[
c_{T}=2.069.
\]

\end{minipage}\par
\par\textbf{Step 4: Normal cutoffs.}\nopagebreak\enspace \par\noindent\begin{minipage}{\linewidth}
\par\textit{Cutoff.} The left cumulative target is $\Phi(c_{N})=0.9750$. It requires a positive $z$ because the target is above $0.5$.\[
\Phi(1.96)=0.9750.
\]
This printed entry matches the target.\[
c_{N}=1.96.
\]
\end{minipage}\par

\par\noindent\begin{minipage}{\linewidth}
\par\textbf{Step 5: Back-transform to contextual units.}\nopagebreak\enspace \[
x=m+ac.
\]
Multiply each selected standardized cutoff by the scale, then add the location. Retain the unrounded result for the comparison errors.
\par\textit{Cutoff back-transformation.}
\[
\begin{aligned}x_T&=3+8(2.069)=3+(16.552)\\&=19.552\approx\boxed{19.55\ \text{minutes}},\\x_N&=3+8(1.96)=3+(15.68)\\&=18.68\approx\boxed{18.68\ \text{minutes}}.\end{aligned}
\]

\end{minipage}\par
\par\textbf{Step 6: Comparison errors.}\nopagebreak\enspace \par\noindent\begin{minipage}{\linewidth}
\par\textit{Cutoff comparison.}
\[
E_A=|r_T-r_N|,\qquad E_{Pt}=\frac{100E_A}{|r_T|},\qquad E_{Pn}=\frac{100E_A}{|r_N|}.
\]
Here $r_T$ and $r_N$ are the unrounded contextual cutoffs being compared.
\[
\begin{aligned}
E_A&=|19.552-(18.68)|=0.872\approx 0.87\ \text{minutes},\\
E_{Pt}&=\frac{100|19.552-(18.68)|}{|19.552|}=\frac{100(0.872)}{19.552}\approx 4.46\%,\\
E_{Pn}&=\frac{100|19.552-(18.68)|}{|18.68|}=\frac{100(0.872)}{18.68}\approx 4.67\%.
\end{aligned}
\]
\end{minipage}\par

\par\noindent\begin{minipage}{\linewidth}
\par\textbf{Step 7: Interpretation.}\nopagebreak\enspace About $2.5\%$ of journeys exceed a delay of $19.55$ minutes under $t$, or $18.68$ minutes under normal. Both positive cutoffs mean late arrival; the normal cutoff is about $0.87$ minutes lower.
\end{minipage}\par
\item \begin{minipage}[t]{\linewidth}\textbf{Electricity forecasting ($\nu=18$).}\nopagebreak
\par\textbf{Step 1: Information and required cutoffs.}\nopagebreak\enspace $n=19$, $m=10$, $a=40$, so $\nu=n-1=19-1=18$. The quantity is electricity-use deviation in kWh. Negative deviations mean use below forecast; positive deviations mean above forecast.\[
P(L<X<U)=0.90\quad\text{with equal excluded tails}.
\]
A central interval splits the excluded probability equally:\[
\alpha=\frac{1-\text{central probability}}{2}=\frac{1-0.90}{2}=\frac{0.10}{2}=0.05.
\]
The lower and upper left cumulative targets are therefore\[
P(X<L)=0.0500,\qquad P(X<U)=1-0.05=0.9500.
\]

\end{minipage}\par
\par\noindent\begin{minipage}{\linewidth}
\par\textbf{Step 2: Standardize the unknowns.}\nopagebreak\enspace \[
c=\frac{x-m}{a}.
\]
Apply the formula separately to the lower and upper unknown limits:
\[
\begin{aligned}c_L&=\frac{L-(10)}{40},\\c_U&=\frac{U-(10)}{40}.\end{aligned}
\]
A negative standardized cutoff lies below the model location $m$; a positive one lies above it. The numerical standardized values are found from the probability targets before back-transformation.
\end{minipage}\par
\par\noindent\begin{minipage}{\linewidth}
\par\textbf{Step 3: Student's $t$ cutoffs.}\nopagebreak\enspace \par\textit{Lower.} Seek $F_{18}(c_{T,L})=0.0500$. This is below $0.5$, so the cutoff is negative. By symmetry use right-tail column $\alpha=0.05$. Row $18$, column $\alpha=0.05$, gives $t^*=1.734$:\[
P(T>1.734)\approx0.0500.
\]
\[
F_{18}(-1.734)=P(T>1.734)\approx0.0500,\qquad c_{T,L}=-1.734.
\]
\par\textit{Upper.} Seek $F_{18}(c_{T,U})=0.9500$. This is above $0.5$, so the cutoff is positive. The right-tail column is $\alpha=1-0.9500=0.05$. Row $18$, column $\alpha=0.05$, gives $t^*=1.734$:\[
P(T>1.734)\approx0.0500.
\]
\[
F_{18}(1.734)=1-P(T>1.734)\approx1-0.0500=0.9500.
\]
\[
c_{T,U}=1.734.
\]

\end{minipage}\par
\par\textbf{Step 4: Normal cutoffs.}\nopagebreak\enspace \par\noindent\begin{minipage}{\linewidth}
\par\textit{Lower.} The left cumulative target is $\Phi(c_{N,L})=0.0500$. It requires a negative $z$ because the target is below $0.5$. For a negative lookup use $\Phi(-z)=1-\Phi(z)$.\[
\Phi(-1.64)=1-0.9495=0.0505.
\]
The entry $\Phi(-1.65)=1-0.9505=0.0495$ is equally close; choose $-1.64$ by the stated tie convention.\[
c_{N,L}=-1.64.
\]
\end{minipage}\par
\par\noindent\begin{minipage}{\linewidth}
\par\textit{Upper.} The left cumulative target is $\Phi(c_{N,U})=0.9500$. It requires a positive $z$ because the target is above $0.5$.\[
\Phi(1.64)=0.9495.
\]
The entry $\Phi(1.65)=0.9505$ is equally close; choose $1.64$ by the stated tie convention.\[
c_{N,U}=1.64.
\]
\end{minipage}\par

\par\noindent\begin{minipage}{\linewidth}
\par\textbf{Step 5: Back-transform to contextual units.}\nopagebreak\enspace \[
x=m+ac.
\]
Multiply each selected standardized cutoff by the scale, then add the location. Retain the unrounded result for the comparison errors.
\par\textit{Lower back-transformation.}
\[
\begin{aligned}L_T&=10+40(-1.734)=10+(-69.36)\\&=-59.36\approx\boxed{-59.36\ \text{kWh}},\\L_N&=10+40(-1.64)=10+(-65.6)\\&=-55.6\approx\boxed{-55.60\ \text{kWh}}.\end{aligned}
\]
\par\textit{Upper back-transformation.}
\[
\begin{aligned}U_T&=10+40(1.734)=10+(69.36)\\&=79.36\approx\boxed{79.36\ \text{kWh}},\\U_N&=10+40(1.64)=10+(65.6)\\&=75.6\approx\boxed{75.60\ \text{kWh}}.\end{aligned}
\]
The two completed intervals, in kWh, are
\[
\boxed{(-59.36,\ 79.36)\quad(t)},\qquad\boxed{(-55.60,\ 75.60)\quad(N)}.
\]

\end{minipage}\par
\par\textbf{Step 6: Comparison errors.}\nopagebreak\enspace \par\noindent\begin{minipage}{\linewidth}
\par\textit{Lower cutoff comparison.}
\[
E_A=|r_T-r_N|,\qquad E_{Pt}=\frac{100E_A}{|r_T|},\qquad E_{Pn}=\frac{100E_A}{|r_N|}.
\]
Here $r_T$ and $r_N$ are the unrounded contextual cutoffs being compared.
\[
\begin{aligned}
E_A&=|-59.36-(-55.6)|=3.76\ \text{kWh},\\
E_{Pt}&=\frac{100|-59.36-(-55.6)|}{|-59.36|}=\frac{100(3.76)}{59.36}\approx 6.33\%,\\
E_{Pn}&=\frac{100|-59.36-(-55.6)|}{|-55.6|}=\frac{100(3.76)}{55.6}\approx 6.76\%.
\end{aligned}
\]
\end{minipage}\par
\par\noindent\begin{minipage}{\linewidth}
\par\textit{Upper cutoff comparison.}
\[
E_A=|r_T-r_N|,\qquad E_{Pt}=\frac{100E_A}{|r_T|},\qquad E_{Pn}=\frac{100E_A}{|r_N|}.
\]
Here $r_T$ and $r_N$ are the unrounded contextual cutoffs being compared.
\[
\begin{aligned}
E_A&=|79.36-(75.6)|=3.76\ \text{kWh},\\
E_{Pt}&=\frac{100|79.36-(75.6)|}{|79.36|}=\frac{100(3.76)}{79.36}\approx 4.74\%,\\
E_{Pn}&=\frac{100|79.36-(75.6)|}{|75.6|}=\frac{100(3.76)}{75.6}\approx 4.97\%.
\end{aligned}
\]
\end{minipage}\par

\par\noindent\begin{minipage}{\linewidth}
\par\textbf{Step 7: Interpretation.}\nopagebreak\enspace The central $90\%$ limits are $(-59.36,79.36)$ kWh under $t$ and $(-55.60,75.60)$ kWh under normal, with approximately $5\%$ in each excluded tail. The lower endpoint is below forecast and the upper endpoint above it. The normal interval extends less far from the mean at both ends, by $3.76$ kWh per endpoint.
\end{minipage}\par
\item \begin{minipage}[t]{\linewidth}\textbf{Packaging inspection ($\nu=5$).}\nopagebreak
\par\textbf{Step 1: Information and required cutoffs.}\nopagebreak\enspace $n=6$, $m=-0.5$, $a=1.5$, so $\nu=n-1=6-1=5$. The quantity is package mass deviation in grams. Negative deviations mean below target mass; positive deviations mean above target.\[
P(X<L)=P(X>U)=0.025.
\]
The central probability is $1-0.05=0.95$. Split the excluded area equally:\[
\alpha=\frac{1-\text{central probability}}{2}=\frac{1-0.95}{2}=\frac{0.05}{2}=0.025.
\]
The left cumulative targets are\[
P(X<L)=0.0250,\qquad P(X<U)=1-0.025=0.9750.
\]

\end{minipage}\par
\par\noindent\begin{minipage}{\linewidth}
\par\textbf{Step 2: Standardize the unknowns.}\nopagebreak\enspace \[
c=\frac{x-m}{a}.
\]
Apply the formula separately to the lower and upper unknown limits:
\[
\begin{aligned}c_L&=\frac{L-(-0.5)}{1.5}=\frac{L+0.5}{1.5},\\c_U&=\frac{U-(-0.5)}{1.5}=\frac{U+0.5}{1.5}.\end{aligned}
\]
A negative standardized cutoff lies below the model location $m$; a positive one lies above it. The numerical standardized values are found from the probability targets before back-transformation.
\end{minipage}\par
\par\noindent\begin{minipage}{\linewidth}
\par\textbf{Step 3: Student's $t$ cutoffs.}\nopagebreak\enspace \par\textit{Lower.} Seek $F_{5}(c_{T,L})=0.0250$. This is below $0.5$, so the cutoff is negative. By symmetry use right-tail column $\alpha=0.025$. Row $5$, column $\alpha=0.025$, gives $t^*=2.571$:\[
P(T>2.571)\approx0.0250.
\]
\[
F_{5}(-2.571)=P(T>2.571)\approx0.0250,\qquad c_{T,L}=-2.571.
\]
\par\textit{Upper.} Seek $F_{5}(c_{T,U})=0.9750$. This is above $0.5$, so the cutoff is positive. The right-tail column is $\alpha=1-0.9750=0.025$. Row $5$, column $\alpha=0.025$, gives $t^*=2.571$:\[
P(T>2.571)\approx0.0250.
\]
\[
F_{5}(2.571)=1-P(T>2.571)\approx1-0.0250=0.9750.
\]
\[
c_{T,U}=2.571.
\]

\end{minipage}\par
\par\textbf{Step 4: Normal cutoffs.}\nopagebreak\enspace \par\noindent\begin{minipage}{\linewidth}
\par\textit{Lower.} The left cumulative target is $\Phi(c_{N,L})=0.0250$. It requires a negative $z$ because the target is below $0.5$. For a negative lookup use $\Phi(-z)=1-\Phi(z)$.\[
\Phi(-1.96)=1-\Phi(1.96)=1-0.9750=0.0250.
\]
This printed entry matches the target.\[
c_{N,L}=-1.96.
\]
\end{minipage}\par
\par\noindent\begin{minipage}{\linewidth}
\par\textit{Upper.} The left cumulative target is $\Phi(c_{N,U})=0.9750$. It requires a positive $z$ because the target is above $0.5$.\[
\Phi(1.96)=0.9750.
\]
This printed entry matches the target.\[
c_{N,U}=1.96.
\]
\end{minipage}\par

\par\noindent\begin{minipage}{\linewidth}
\par\textbf{Step 5: Back-transform to contextual units.}\nopagebreak\enspace \[
x=m+ac.
\]
Multiply each selected standardized cutoff by the scale, then add the location. Retain the unrounded result for the comparison errors.
\par\textit{Lower back-transformation.}
\[
\begin{aligned}L_T&=-0.5+1.5(-2.571)=-0.5+(-3.8565)\\&=-4.3565\approx\boxed{-4.36\ \text{g}},\\L_N&=-0.5+1.5(-1.96)=-0.5+(-2.94)\\&=-3.44\approx\boxed{-3.44\ \text{g}}.\end{aligned}
\]
\par\textit{Upper back-transformation.}
\[
\begin{aligned}U_T&=-0.5+1.5(2.571)=-0.5+(3.8565)\\&=3.3565\approx\boxed{3.36\ \text{g}},\\U_N&=-0.5+1.5(1.96)=-0.5+(2.94)\\&=2.44\approx\boxed{2.44\ \text{g}}.\end{aligned}
\]
The two completed intervals, in g, are
\[
\boxed{(-4.36,\ 3.36)\quad(t)},\qquad\boxed{(-3.44,\ 2.44)\quad(N)}.
\]

\end{minipage}\par
\par\textbf{Step 6: Comparison errors.}\nopagebreak\enspace \par\noindent\begin{minipage}{\linewidth}
\par\textit{Lower cutoff comparison.}
\[
E_A=|r_T-r_N|,\qquad E_{Pt}=\frac{100E_A}{|r_T|},\qquad E_{Pn}=\frac{100E_A}{|r_N|}.
\]
Here $r_T$ and $r_N$ are the unrounded contextual cutoffs being compared.
\[
\begin{aligned}
E_A&=|-4.3565-(-3.44)|=0.9165\approx 0.92\ \text{g},\\
E_{Pt}&=\frac{100|-4.3565-(-3.44)|}{|-4.3565|}=\frac{100(0.9165)}{4.3565}\approx 21.04\%,\\
E_{Pn}&=\frac{100|-4.3565-(-3.44)|}{|-3.44|}=\frac{100(0.9165)}{3.44}\approx 26.64\%.
\end{aligned}
\]
\end{minipage}\par
\par\noindent\begin{minipage}{\linewidth}
\par\textit{Upper cutoff comparison.}
\[
E_A=|r_T-r_N|,\qquad E_{Pt}=\frac{100E_A}{|r_T|},\qquad E_{Pn}=\frac{100E_A}{|r_N|}.
\]
Here $r_T$ and $r_N$ are the unrounded contextual cutoffs being compared.
\[
\begin{aligned}
E_A&=|3.3565-(2.44)|=0.9165\approx 0.92\ \text{g},\\
E_{Pt}&=\frac{100|3.3565-(2.44)|}{|3.3565|}=\frac{100(0.9165)}{3.3565}\approx 27.31\%,\\
E_{Pn}&=\frac{100|3.3565-(2.44)|}{|2.44|}=\frac{100(0.9165)}{2.44}\approx 37.56\%.
\end{aligned}
\]
\end{minipage}\par

\par\noindent\begin{minipage}{\linewidth}
\par\textbf{Step 7: Interpretation.}\nopagebreak\enspace The limits are $(-4.36,3.36)$ g under $t$ and $(-3.44,2.44)$ g under normal, enclosing approximately $95\%$ of mass deviations with $2.5\%$ in each tail. The lower limits are below target and the upper limits above target. Each normal endpoint is about $0.92$ g closer to the mean; the percentage errors differ because the endpoint magnitudes differ.
\end{minipage}\par
\item \begin{minipage}[t]{\linewidth}\textbf{Processing-time alerts ($\nu=12$).}\nopagebreak
\par\textbf{Step 1: Information and required cutoffs.}\nopagebreak\enspace $n=13$, $m=1$, $a=5$, so $\nu=n-1=13-1=12$. The quantity is processing-time deviation in minutes. Negative deviations mean processing faster than target; positive deviations mean slower.\[
P(X<L)=0.05,\qquad P(X>U)=0.10.
\]
The two tails are unequal, so use their specified probabilities separately. Do not split $0.15$ into equal tails. The lower target is already cumulative; complement the upper right tail:\[
P(X<L)=0.0500,\qquad P(X<U)=1-0.10=0.9000.
\]
The area between the limits is $1-0.05-0.10=0.85$.
\end{minipage}\par
\par\noindent\begin{minipage}{\linewidth}
\par\textbf{Step 2: Standardize the unknowns.}\nopagebreak\enspace \[
c=\frac{x-m}{a}.
\]
Apply the formula separately to the lower and upper unknown limits:
\[
\begin{aligned}c_L&=\frac{L-(1)}{5},\\c_U&=\frac{U-(1)}{5}.\end{aligned}
\]
A negative standardized cutoff lies below the model location $m$; a positive one lies above it. The numerical standardized values are found from the probability targets before back-transformation.
\end{minipage}\par
\par\noindent\begin{minipage}{\linewidth}
\par\textbf{Step 3: Student's $t$ cutoffs.}\nopagebreak\enspace \par\textit{Lower.} Seek $F_{12}(c_{T,L})=0.0500$. This is below $0.5$, so the cutoff is negative. By symmetry use right-tail column $\alpha=0.05$. Row $12$, column $\alpha=0.05$, gives $t^*=1.782$:\[
P(T>1.782)\approx0.0500.
\]
\[
F_{12}(-1.782)=P(T>1.782)\approx0.0500,\qquad c_{T,L}=-1.782.
\]
\par\textit{Upper.} Seek $F_{12}(c_{T,U})=0.9000$. This is above $0.5$, so the cutoff is positive. The right-tail column is $\alpha=1-0.9000=0.1$. Row $12$, column $\alpha=0.1$, gives $t^*=1.356$:\[
P(T>1.356)\approx0.1000.
\]
\[
F_{12}(1.356)=1-P(T>1.356)\approx1-0.1000=0.9000.
\]
\[
c_{T,U}=1.356.
\]

\end{minipage}\par
\par\textbf{Step 4: Normal cutoffs.}\nopagebreak\enspace \par\noindent\begin{minipage}{\linewidth}
\par\textit{Lower.} The left cumulative target is $\Phi(c_{N,L})=0.0500$. It requires a negative $z$ because the target is below $0.5$. For a negative lookup use $\Phi(-z)=1-\Phi(z)$.\[
\Phi(-1.64)=1-0.9495=0.0505.
\]
The entry $\Phi(-1.65)=1-0.9505=0.0495$ is equally close; choose $-1.64$ by the stated tie convention.\[
c_{N,L}=-1.64.
\]
\end{minipage}\par
\par\noindent\begin{minipage}{\linewidth}
\par\textit{Upper.} The left cumulative target is $\Phi(c_{N,U})=0.9000$. It requires a positive $z$ because the target is above $0.5$.\[
\Phi(1.28)=0.8997.
\]
The adjacent entry $\Phi(1.29)=0.9015$ is farther from $0.9000$.\[
c_{N,U}=1.28.
\]
\end{minipage}\par

\par\noindent\begin{minipage}{\linewidth}
\par\textbf{Step 5: Back-transform to contextual units.}\nopagebreak\enspace \[
x=m+ac.
\]
Multiply each selected standardized cutoff by the scale, then add the location. Retain the unrounded result for the comparison errors.
\par\textit{Lower back-transformation.}
\[
\begin{aligned}L_T&=1+5(-1.782)=1+(-8.91)\\&=-7.91\approx\boxed{-7.91\ \text{minutes}},\\L_N&=1+5(-1.64)=1+(-8.2)\\&=-7.2\approx\boxed{-7.20\ \text{minutes}}.\end{aligned}
\]
\par\textit{Upper back-transformation.}
\[
\begin{aligned}U_T&=1+5(1.356)=1+(6.78)\\&=7.78\approx\boxed{7.78\ \text{minutes}},\\U_N&=1+5(1.28)=1+(6.4)\\&=7.4\approx\boxed{7.40\ \text{minutes}}.\end{aligned}
\]
The two completed intervals, in minutes, are
\[
\boxed{(-7.91,\ 7.78)\quad(t)},\qquad\boxed{(-7.20,\ 7.40)\quad(N)}.
\]

\end{minipage}\par
\par\textbf{Step 6: Comparison errors.}\nopagebreak\enspace \par\noindent\begin{minipage}{\linewidth}
\par\textit{Lower cutoff comparison.}
\[
E_A=|r_T-r_N|,\qquad E_{Pt}=\frac{100E_A}{|r_T|},\qquad E_{Pn}=\frac{100E_A}{|r_N|}.
\]
Here $r_T$ and $r_N$ are the unrounded contextual cutoffs being compared.
\[
\begin{aligned}
E_A&=|-7.91-(-7.2)|=0.71\ \text{minutes},\\
E_{Pt}&=\frac{100|-7.91-(-7.2)|}{|-7.91|}=\frac{100(0.71)}{7.91}\approx 8.98\%,\\
E_{Pn}&=\frac{100|-7.91-(-7.2)|}{|-7.2|}=\frac{100(0.71)}{7.2}\approx 9.86\%.
\end{aligned}
\]
\end{minipage}\par
\par\noindent\begin{minipage}{\linewidth}
\par\textit{Upper cutoff comparison.}
\[
E_A=|r_T-r_N|,\qquad E_{Pt}=\frac{100E_A}{|r_T|},\qquad E_{Pn}=\frac{100E_A}{|r_N|}.
\]
Here $r_T$ and $r_N$ are the unrounded contextual cutoffs being compared.
\[
\begin{aligned}
E_A&=|7.78-(7.4)|=0.38\ \text{minutes},\\
E_{Pt}&=\frac{100|7.78-(7.4)|}{|7.78|}=\frac{100(0.38)}{7.78}\approx 4.88\%,\\
E_{Pn}&=\frac{100|7.78-(7.4)|}{|7.4|}=\frac{100(0.38)}{7.4}\approx 5.14\%.
\end{aligned}
\]
\end{minipage}\par

\par\noindent\begin{minipage}{\linewidth}
\par\textbf{Step 7: Interpretation.}\nopagebreak\enspace The alert limits are $(-7.91,7.78)$ minutes under $t$ and $(-7.20,7.40)$ minutes under normal. Each model leaves approximately $5\%$ below its lower limit, $10\%$ above its upper limit, and $85\%$ between them. The negative lower limits mean faster processing and positive upper limits mean slower processing. Normal moves the lower endpoint inward by $0.71$ minutes and the upper endpoint inward by $0.38$ minutes; the unequal tails give unequal distances from the mean.
\end{minipage}\par
\item \begin{minipage}[t]{\linewidth}\textbf{Hospital waiting ($\nu=12$).}\nopagebreak
\par\textbf{Step 1: Information and required cutoffs.}\nopagebreak\enspace $n=13$, $m=1$, $a=5$, so $\nu=n-1=13-1=12$. The quantity is waiting-time deviation from schedule in minutes. Negative deviations mean earlier than scheduled.\[
P(X<x)=0.10.
\]
The requested probability is already a left cumulative target: $P(X<x)=0.1000$. The remaining area is $P(X>x)=1-0.10=0.90$.
\end{minipage}\par
\par\noindent\begin{minipage}{\linewidth}
\par\textbf{Step 2: Standardize the unknowns.}\nopagebreak\enspace \[
c=\frac{x-m}{a}.
\]
\[
c=\frac{x-(1)}{5}.
\]
A negative standardized cutoff lies below the model location $m$; a positive one lies above it. The numerical standardized values are found from the probability targets before back-transformation.
\end{minipage}\par
\par\noindent\begin{minipage}{\linewidth}
\par\textbf{Step 3: Student's $t$ cutoffs.}\nopagebreak\enspace \par\textit{Cutoff.} Seek $F_{12}(c_{T})=0.1000$. This is below $0.5$, so the cutoff is negative. By symmetry use right-tail column $\alpha=0.1$. Row $12$, column $\alpha=0.1$, gives $t^*=1.356$:\[
P(T>1.356)\approx0.1000.
\]
\[
F_{12}(-1.356)=P(T>1.356)\approx0.1000,\qquad c_{T}=-1.356.
\]

\end{minipage}\par
\par\textbf{Step 4: Normal cutoffs.}\nopagebreak\enspace \par\noindent\begin{minipage}{\linewidth}
\par\textit{Cutoff.} The left cumulative target is $\Phi(c_{N})=0.1000$. It requires a negative $z$ because the target is below $0.5$. For a negative lookup use $\Phi(-z)=1-\Phi(z)$.\[
\Phi(-1.28)=1-\Phi(1.28)=1-0.8997=0.1003.
\]
The adjacent entry $\Phi(-1.29)=1-0.9015=0.0985$ is farther from $0.1000$.\[
c_{N}=-1.28.
\]
\end{minipage}\par

\par\noindent\begin{minipage}{\linewidth}
\par\textbf{Step 5: Back-transform to contextual units.}\nopagebreak\enspace \[
x=m+ac.
\]
Multiply each selected standardized cutoff by the scale, then add the location. Retain the unrounded result for the comparison errors.
\par\textit{Cutoff back-transformation.}
\[
\begin{aligned}x_T&=1+5(-1.356)=1+(-6.78)\\&=-5.78\approx\boxed{-5.78\ \text{minutes}},\\x_N&=1+5(-1.28)=1+(-6.4)\\&=-5.4\approx\boxed{-5.40\ \text{minutes}}.\end{aligned}
\]

\end{minipage}\par
\par\textbf{Step 6: Comparison errors.}\nopagebreak\enspace \par\noindent\begin{minipage}{\linewidth}
\par\textit{Cutoff comparison.}
\[
E_A=|r_T-r_N|,\qquad E_{Pt}=\frac{100E_A}{|r_T|},\qquad E_{Pn}=\frac{100E_A}{|r_N|}.
\]
Here $r_T$ and $r_N$ are the unrounded contextual cutoffs being compared.
\[
\begin{aligned}
E_A&=|-5.78-(-5.4)|=0.38\ \text{minutes},\\
E_{Pt}&=\frac{100|-5.78-(-5.4)|}{|-5.78|}=\frac{100(0.38)}{5.78}\approx 6.57\%,\\
E_{Pn}&=\frac{100|-5.78-(-5.4)|}{|-5.4|}=\frac{100(0.38)}{5.4}\approx 7.04\%.
\end{aligned}
\]
\end{minipage}\par

\par\noindent\begin{minipage}{\linewidth}
\par\textbf{Step 7: Interpretation.}\nopagebreak\enspace About $10\%$ of waiting-time deviations lie below $-5.78$ minutes under $t$, or below $-5.40$ minutes under normal. Thus the thresholds represent about $5.78$ or $5.40$ minutes earlier than scheduled, respectively. The normal cutoff is $0.38$ minutes closer to the mean.
\end{minipage}\par
\item \begin{minipage}[t]{\linewidth}\textbf{Production tolerance ($\nu=20$).}\nopagebreak
\par\textbf{Step 1: Information and required cutoffs.}\nopagebreak\enspace $n=21$, $m=0$, $a=12$, so $\nu=n-1=21-1=20$. The quantity is production-output deviation in units. Negative deviations mean below target output; positive deviations mean above target.\[
P(L<X<U)=0.95\quad\text{with equal excluded tails}.
\]
Calculate the area in each tail first:\[
\alpha=\frac{1-\text{central probability}}{2}=\frac{1-0.95}{2}=\frac{0.05}{2}=0.025.
\]
The corresponding left cumulative targets are\[
P(X<L)=0.0250,\qquad P(X<U)=1-0.025=0.9750.
\]

\end{minipage}\par
\par\noindent\begin{minipage}{\linewidth}
\par\textbf{Step 2: Standardize the unknowns.}\nopagebreak\enspace \[
c=\frac{x-m}{a}.
\]
Apply the formula separately to the lower and upper unknown limits:
\[
\begin{aligned}c_L&=\frac{L-(0)}{12},\\c_U&=\frac{U-(0)}{12}.\end{aligned}
\]
A negative standardized cutoff lies below the model location $m$; a positive one lies above it. The numerical standardized values are found from the probability targets before back-transformation.
\end{minipage}\par
\par\noindent\begin{minipage}{\linewidth}
\par\textbf{Step 3: Student's $t$ cutoffs.}\nopagebreak\enspace \par\textit{Lower.} Seek $F_{20}(c_{T,L})=0.0250$. This is below $0.5$, so the cutoff is negative. By symmetry use right-tail column $\alpha=0.025$. Row $20$, column $\alpha=0.025$, gives $t^*=2.086$:\[
P(T>2.086)\approx0.0250.
\]
\[
F_{20}(-2.086)=P(T>2.086)\approx0.0250,\qquad c_{T,L}=-2.086.
\]
\par\textit{Upper.} Seek $F_{20}(c_{T,U})=0.9750$. This is above $0.5$, so the cutoff is positive. The right-tail column is $\alpha=1-0.9750=0.025$. Row $20$, column $\alpha=0.025$, gives $t^*=2.086$:\[
P(T>2.086)\approx0.0250.
\]
\[
F_{20}(2.086)=1-P(T>2.086)\approx1-0.0250=0.9750.
\]
\[
c_{T,U}=2.086.
\]

\end{minipage}\par
\par\textbf{Step 4: Normal cutoffs.}\nopagebreak\enspace \par\noindent\begin{minipage}{\linewidth}
\par\textit{Lower.} The left cumulative target is $\Phi(c_{N,L})=0.0250$. It requires a negative $z$ because the target is below $0.5$. For a negative lookup use $\Phi(-z)=1-\Phi(z)$.\[
\Phi(-1.96)=1-\Phi(1.96)=1-0.9750=0.0250.
\]
This printed entry matches the target.\[
c_{N,L}=-1.96.
\]
\end{minipage}\par
\par\noindent\begin{minipage}{\linewidth}
\par\textit{Upper.} The left cumulative target is $\Phi(c_{N,U})=0.9750$. It requires a positive $z$ because the target is above $0.5$.\[
\Phi(1.96)=0.9750.
\]
This printed entry matches the target.\[
c_{N,U}=1.96.
\]
\end{minipage}\par

\par\noindent\begin{minipage}{\linewidth}
\par\textbf{Step 5: Back-transform to contextual units.}\nopagebreak\enspace \[
x=m+ac.
\]
Multiply each selected standardized cutoff by the scale, then add the location. Retain the unrounded result for the comparison errors.
\par\textit{Lower back-transformation.}
\[
\begin{aligned}L_T&=0+12(-2.086)=0+(-25.032)\\&=-25.032\approx\boxed{-25.03\ \text{units}},\\L_N&=0+12(-1.96)=0+(-23.52)\\&=-23.52\approx\boxed{-23.52\ \text{units}}.\end{aligned}
\]
\par\textit{Upper back-transformation.}
\[
\begin{aligned}U_T&=0+12(2.086)=0+(25.032)\\&=25.032\approx\boxed{25.03\ \text{units}},\\U_N&=0+12(1.96)=0+(23.52)\\&=23.52\approx\boxed{23.52\ \text{units}}.\end{aligned}
\]
The two completed intervals, in units, are
\[
\boxed{(-25.03,\ 25.03)\quad(t)},\qquad\boxed{(-23.52,\ 23.52)\quad(N)}.
\]

\end{minipage}\par
\par\textbf{Step 6: Comparison errors.}\nopagebreak\enspace \par\noindent\begin{minipage}{\linewidth}
\par\textit{Lower cutoff comparison.}
\[
E_A=|r_T-r_N|,\qquad E_{Pt}=\frac{100E_A}{|r_T|},\qquad E_{Pn}=\frac{100E_A}{|r_N|}.
\]
Here $r_T$ and $r_N$ are the unrounded contextual cutoffs being compared.
\[
\begin{aligned}
E_A&=|-25.032-(-23.52)|=1.512\approx 1.51\ \text{units},\\
E_{Pt}&=\frac{100|-25.032-(-23.52)|}{|-25.032|}=\frac{100(1.512)}{25.032}\approx 6.04\%,\\
E_{Pn}&=\frac{100|-25.032-(-23.52)|}{|-23.52|}=\frac{100(1.512)}{23.52}\approx 6.43\%.
\end{aligned}
\]
\end{minipage}\par
\par\noindent\begin{minipage}{\linewidth}
\par\textit{Upper cutoff comparison.}
\[
E_A=|r_T-r_N|,\qquad E_{Pt}=\frac{100E_A}{|r_T|},\qquad E_{Pn}=\frac{100E_A}{|r_N|}.
\]
Here $r_T$ and $r_N$ are the unrounded contextual cutoffs being compared.
\[
\begin{aligned}
E_A&=|25.032-(23.52)|=1.512\approx 1.51\ \text{units},\\
E_{Pt}&=\frac{100|25.032-(23.52)|}{|25.032|}=\frac{100(1.512)}{25.032}\approx 6.04\%,\\
E_{Pn}&=\frac{100|25.032-(23.52)|}{|23.52|}=\frac{100(1.512)}{23.52}\approx 6.43\%.
\end{aligned}
\]
\end{minipage}\par

\par\noindent\begin{minipage}{\linewidth}
\par\textbf{Step 7: Interpretation.}\nopagebreak\enspace The central $95\%$ limits are $(-25.03,25.03)$ units under $t$ and $(-23.52,23.52)$ units under normal, leaving approximately $2.5\%$ in each tail. They describe output below and above target, symmetrically because $m=0$. Each normal endpoint is about $1.51$ units closer to the target.
\end{minipage}\par
\item \begin{minipage}[t]{\linewidth}\textbf{Electricity peak ($\nu=15$).}\nopagebreak
\par\textbf{Step 1: Information and required cutoffs.}\nopagebreak\enspace $n=16$, $m=3$, $a=7$, so $\nu=n-1=16-1=15$. The quantity is peak-demand deviation in MW. A positive deviation means peak demand above forecast.\[
P(X>x)=0.01.
\]
Convert the right-tail condition to the left cumulative target:\[
P(X<x)=1-P(X>x)=1-0.01=0.9900.
\]

\end{minipage}\par
\par\noindent\begin{minipage}{\linewidth}
\par\textbf{Step 2: Standardize the unknowns.}\nopagebreak\enspace \[
c=\frac{x-m}{a}.
\]
\[
c=\frac{x-(3)}{7}.
\]
A negative standardized cutoff lies below the model location $m$; a positive one lies above it. The numerical standardized values are found from the probability targets before back-transformation.
\end{minipage}\par
\par\noindent\begin{minipage}{\linewidth}
\par\textbf{Step 3: Student's $t$ cutoffs.}\nopagebreak\enspace \par\textit{Cutoff.} Seek $F_{15}(c_{T})=0.9900$. This is above $0.5$, so the cutoff is positive. The right-tail column is $\alpha=1-0.9900=0.01$. Row $15$, column $\alpha=0.01$, gives $t^*=2.602$:\[
P(T>2.602)\approx0.0100.
\]
\[
F_{15}(2.602)=1-P(T>2.602)\approx1-0.0100=0.9900.
\]
\[
c_{T}=2.602.
\]

\end{minipage}\par
\par\textbf{Step 4: Normal cutoffs.}\nopagebreak\enspace \par\noindent\begin{minipage}{\linewidth}
\par\textit{Cutoff.} The left cumulative target is $\Phi(c_{N})=0.9900$. It requires a positive $z$ because the target is above $0.5$.\[
\Phi(2.33)=0.9901.
\]
The adjacent entry $\Phi(2.32)=0.9898$ is farther from $0.9900$.\[
c_{N}=2.33.
\]
\end{minipage}\par

\par\noindent\begin{minipage}{\linewidth}
\par\textbf{Step 5: Back-transform to contextual units.}\nopagebreak\enspace \[
x=m+ac.
\]
Multiply each selected standardized cutoff by the scale, then add the location. Retain the unrounded result for the comparison errors.
\par\textit{Cutoff back-transformation.}
\[
\begin{aligned}x_T&=3+7(2.602)=3+(18.214)\\&=21.214\approx\boxed{21.21\ \text{MW}},\\x_N&=3+7(2.33)=3+(16.31)\\&=19.31\approx\boxed{19.31\ \text{MW}}.\end{aligned}
\]

\end{minipage}\par
\par\textbf{Step 6: Comparison errors.}\nopagebreak\enspace \par\noindent\begin{minipage}{\linewidth}
\par\textit{Cutoff comparison.}
\[
E_A=|r_T-r_N|,\qquad E_{Pt}=\frac{100E_A}{|r_T|},\qquad E_{Pn}=\frac{100E_A}{|r_N|}.
\]
Here $r_T$ and $r_N$ are the unrounded contextual cutoffs being compared.
\[
\begin{aligned}
E_A&=|21.214-(19.31)|=1.904\approx 1.90\ \text{MW},\\
E_{Pt}&=\frac{100|21.214-(19.31)|}{|21.214|}=\frac{100(1.904)}{21.214}\approx 8.98\%,\\
E_{Pn}&=\frac{100|21.214-(19.31)|}{|19.31|}=\frac{100(1.904)}{19.31}\approx 9.86\%.
\end{aligned}
\]
\end{minipage}\par

\par\noindent\begin{minipage}{\linewidth}
\par\textbf{Step 7: Interpretation.}\nopagebreak\enspace About $1\%$ of peak-demand deviations exceed $21.21$ MW above forecast under $t$, or $19.31$ MW above forecast under normal. The normal cutoff is about $1.90$ MW lower; each cutoff refers to its own model probability.
\end{minipage}\par
\item \begin{minipage}[t]{\linewidth}\textbf{Transit reliability ($\nu=8$).}\nopagebreak
\par\textbf{Step 1: Information and required cutoffs.}\nopagebreak\enspace $n=9$, $m=-1$, $a=3$, so $\nu=n-1=9-1=8$. The quantity is arrival-time deviation in minutes. Negative deviations mean early arrival; positive deviations mean late arrival.\[
P(X<L)=0.05,\qquad P(X>U)=0.025.
\]
The alert tails are unequal: use $0.05$ for the lower tail and $0.025$ for the upper tail separately. The left cumulative targets are\[
P(X<L)=0.0500,\qquad P(X<U)=1-0.025=0.9750.
\]
The enclosed area is $1-0.05-0.025=0.925$, not a central interval with equal tails.
\end{minipage}\par
\par\noindent\begin{minipage}{\linewidth}
\par\textbf{Step 2: Standardize the unknowns.}\nopagebreak\enspace \[
c=\frac{x-m}{a}.
\]
Apply the formula separately to the lower and upper unknown limits:
\[
\begin{aligned}c_L&=\frac{L-(-1)}{3}=\frac{L+1}{3},\\c_U&=\frac{U-(-1)}{3}=\frac{U+1}{3}.\end{aligned}
\]
A negative standardized cutoff lies below the model location $m$; a positive one lies above it. The numerical standardized values are found from the probability targets before back-transformation.
\end{minipage}\par
\par\noindent\begin{minipage}{\linewidth}
\par\textbf{Step 3: Student's $t$ cutoffs.}\nopagebreak\enspace \par\textit{Lower.} Seek $F_{8}(c_{T,L})=0.0500$. This is below $0.5$, so the cutoff is negative. By symmetry use right-tail column $\alpha=0.05$. Row $8$, column $\alpha=0.05$, gives $t^*=1.860$:\[
P(T>1.860)\approx0.0500.
\]
\[
F_{8}(-1.860)=P(T>1.860)\approx0.0500,\qquad c_{T,L}=-1.860.
\]
\par\textit{Upper.} Seek $F_{8}(c_{T,U})=0.9750$. This is above $0.5$, so the cutoff is positive. The right-tail column is $\alpha=1-0.9750=0.025$. Row $8$, column $\alpha=0.025$, gives $t^*=2.306$:\[
P(T>2.306)\approx0.0250.
\]
\[
F_{8}(2.306)=1-P(T>2.306)\approx1-0.0250=0.9750.
\]
\[
c_{T,U}=2.306.
\]

\end{minipage}\par
\par\textbf{Step 4: Normal cutoffs.}\nopagebreak\enspace \par\noindent\begin{minipage}{\linewidth}
\par\textit{Lower.} The left cumulative target is $\Phi(c_{N,L})=0.0500$. It requires a negative $z$ because the target is below $0.5$. For a negative lookup use $\Phi(-z)=1-\Phi(z)$.\[
\Phi(-1.64)=1-0.9495=0.0505.
\]
The entry $\Phi(-1.65)=1-0.9505=0.0495$ is equally close; choose $-1.64$ by the stated tie convention.\[
c_{N,L}=-1.64.
\]
\end{minipage}\par
\par\noindent\begin{minipage}{\linewidth}
\par\textit{Upper.} The left cumulative target is $\Phi(c_{N,U})=0.9750$. It requires a positive $z$ because the target is above $0.5$.\[
\Phi(1.96)=0.9750.
\]
This printed entry matches the target.\[
c_{N,U}=1.96.
\]
\end{minipage}\par

\par\noindent\begin{minipage}{\linewidth}
\par\textbf{Step 5: Back-transform to contextual units.}\nopagebreak\enspace \[
x=m+ac.
\]
Multiply each selected standardized cutoff by the scale, then add the location. Retain the unrounded result for the comparison errors.
\par\textit{Lower back-transformation.}
\[
\begin{aligned}L_T&=-1+3(-1.860)=-1+(-5.58)\\&=-6.58\approx\boxed{-6.58\ \text{minutes}},\\L_N&=-1+3(-1.64)=-1+(-4.92)\\&=-5.92\approx\boxed{-5.92\ \text{minutes}}.\end{aligned}
\]
\par\textit{Upper back-transformation.}
\[
\begin{aligned}U_T&=-1+3(2.306)=-1+(6.918)\\&=5.918\approx\boxed{5.92\ \text{minutes}},\\U_N&=-1+3(1.96)=-1+(5.88)\\&=4.88\approx\boxed{4.88\ \text{minutes}}.\end{aligned}
\]
The two completed intervals, in minutes, are
\[
\boxed{(-6.58,\ 5.92)\quad(t)},\qquad\boxed{(-5.92,\ 4.88)\quad(N)}.
\]

\end{minipage}\par
\par\textbf{Step 6: Comparison errors.}\nopagebreak\enspace \par\noindent\begin{minipage}{\linewidth}
\par\textit{Lower cutoff comparison.}
\[
E_A=|r_T-r_N|,\qquad E_{Pt}=\frac{100E_A}{|r_T|},\qquad E_{Pn}=\frac{100E_A}{|r_N|}.
\]
Here $r_T$ and $r_N$ are the unrounded contextual cutoffs being compared.
\[
\begin{aligned}
E_A&=|-6.58-(-5.92)|=0.66\ \text{minutes},\\
E_{Pt}&=\frac{100|-6.58-(-5.92)|}{|-6.58|}=\frac{100(0.66)}{6.58}\approx 10.03\%,\\
E_{Pn}&=\frac{100|-6.58-(-5.92)|}{|-5.92|}=\frac{100(0.66)}{5.92}\approx 11.15\%.
\end{aligned}
\]
\end{minipage}\par
\par\noindent\begin{minipage}{\linewidth}
\par\textit{Upper cutoff comparison.}
\[
E_A=|r_T-r_N|,\qquad E_{Pt}=\frac{100E_A}{|r_T|},\qquad E_{Pn}=\frac{100E_A}{|r_N|}.
\]
Here $r_T$ and $r_N$ are the unrounded contextual cutoffs being compared.
\[
\begin{aligned}
E_A&=|5.918-(4.88)|=1.038\approx 1.04\ \text{minutes},\\
E_{Pt}&=\frac{100|5.918-(4.88)|}{|5.918|}=\frac{100(1.038)}{5.918}\approx 17.54\%,\\
E_{Pn}&=\frac{100|5.918-(4.88)|}{|4.88|}=\frac{100(1.038)}{4.88}\approx 21.27\%.
\end{aligned}
\]
\end{minipage}\par

\par\noindent\begin{minipage}{\linewidth}
\par\textbf{Step 7: Interpretation.}\nopagebreak\enspace The alert limits are $(-6.58,5.92)$ minutes under $t$ and $(-5.92,4.88)$ minutes under normal. Approximately $5\%$ of journeys are earlier than the lower limit and $2.5\%$ are later than the upper limit, leaving $92.5\%$ between each model's limits. Normal moves the early limit toward the mean by $0.66$ minutes and the late limit toward the mean by about $1.04$ minutes.
\end{minipage}\par
\item \begin{minipage}[t]{\linewidth}\textbf{Shipping weight ($\nu=10$).}\nopagebreak
\par\textbf{Step 1: Information and required cutoffs.}\nopagebreak\enspace $n=11$, $m=2$, $a=4$, so $\nu=n-1=11-1=10$. The quantity is shipment-mass deviation in kg. Negative deviations mean below target mass; positive deviations mean above target.\[
P(L<X<U)=0.98\quad\text{with equal excluded tails}.
\]
Calculate each excluded tail:\[
\alpha=\frac{1-\text{central probability}}{2}=\frac{1-0.98}{2}=\frac{0.02}{2}=0.01.
\]
The left cumulative targets are\[
P(X<L)=0.0100,\qquad P(X<U)=1-0.01=0.9900.
\]

\end{minipage}\par
\par\noindent\begin{minipage}{\linewidth}
\par\textbf{Step 2: Standardize the unknowns.}\nopagebreak\enspace \[
c=\frac{x-m}{a}.
\]
Apply the formula separately to the lower and upper unknown limits:
\[
\begin{aligned}c_L&=\frac{L-(2)}{4},\\c_U&=\frac{U-(2)}{4}.\end{aligned}
\]
A negative standardized cutoff lies below the model location $m$; a positive one lies above it. The numerical standardized values are found from the probability targets before back-transformation.
\end{minipage}\par
\par\noindent\begin{minipage}{\linewidth}
\par\textbf{Step 3: Student's $t$ cutoffs.}\nopagebreak\enspace \par\textit{Lower.} Seek $F_{10}(c_{T,L})=0.0100$. This is below $0.5$, so the cutoff is negative. By symmetry use right-tail column $\alpha=0.01$. Row $10$, column $\alpha=0.01$, gives $t^*=2.764$:\[
P(T>2.764)\approx0.0100.
\]
\[
F_{10}(-2.764)=P(T>2.764)\approx0.0100,\qquad c_{T,L}=-2.764.
\]
\par\textit{Upper.} Seek $F_{10}(c_{T,U})=0.9900$. This is above $0.5$, so the cutoff is positive. The right-tail column is $\alpha=1-0.9900=0.01$. Row $10$, column $\alpha=0.01$, gives $t^*=2.764$:\[
P(T>2.764)\approx0.0100.
\]
\[
F_{10}(2.764)=1-P(T>2.764)\approx1-0.0100=0.9900.
\]
\[
c_{T,U}=2.764.
\]

\end{minipage}\par
\par\textbf{Step 4: Normal cutoffs.}\nopagebreak\enspace \par\noindent\begin{minipage}{\linewidth}
\par\textit{Lower.} The left cumulative target is $\Phi(c_{N,L})=0.0100$. It requires a negative $z$ because the target is below $0.5$. For a negative lookup use $\Phi(-z)=1-\Phi(z)$.\[
\Phi(-2.33)=1-\Phi(2.33)=1-0.9901=0.0099.
\]
The adjacent entry $\Phi(-2.32)=1-0.9898=0.0102$ is farther from $0.0100$.\[
c_{N,L}=-2.33.
\]
\end{minipage}\par
\par\noindent\begin{minipage}{\linewidth}
\par\textit{Upper.} The left cumulative target is $\Phi(c_{N,U})=0.9900$. It requires a positive $z$ because the target is above $0.5$.\[
\Phi(2.33)=0.9901.
\]
The adjacent entry $\Phi(2.32)=0.9898$ is farther from $0.9900$.\[
c_{N,U}=2.33.
\]
\end{minipage}\par

\par\noindent\begin{minipage}{\linewidth}
\par\textbf{Step 5: Back-transform to contextual units.}\nopagebreak\enspace \[
x=m+ac.
\]
Multiply each selected standardized cutoff by the scale, then add the location. Retain the unrounded result for the comparison errors.
\par\textit{Lower back-transformation.}
\[
\begin{aligned}L_T&=2+4(-2.764)=2+(-11.056)\\&=-9.056\approx\boxed{-9.06\ \text{kg}},\\L_N&=2+4(-2.33)=2+(-9.32)\\&=-7.32\approx\boxed{-7.32\ \text{kg}}.\end{aligned}
\]
\par\textit{Upper back-transformation.}
\[
\begin{aligned}U_T&=2+4(2.764)=2+(11.056)\\&=13.056\approx\boxed{13.06\ \text{kg}},\\U_N&=2+4(2.33)=2+(9.32)\\&=11.32\approx\boxed{11.32\ \text{kg}}.\end{aligned}
\]
The two completed intervals, in kg, are
\[
\boxed{(-9.06,\ 13.06)\quad(t)},\qquad\boxed{(-7.32,\ 11.32)\quad(N)}.
\]

\end{minipage}\par
\par\textbf{Step 6: Comparison errors.}\nopagebreak\enspace \par\noindent\begin{minipage}{\linewidth}
\par\textit{Lower cutoff comparison.}
\[
E_A=|r_T-r_N|,\qquad E_{Pt}=\frac{100E_A}{|r_T|},\qquad E_{Pn}=\frac{100E_A}{|r_N|}.
\]
Here $r_T$ and $r_N$ are the unrounded contextual cutoffs being compared.
\[
\begin{aligned}
E_A&=|-9.056-(-7.32)|=1.736\approx 1.74\ \text{kg},\\
E_{Pt}&=\frac{100|-9.056-(-7.32)|}{|-9.056|}=\frac{100(1.736)}{9.056}\approx 19.17\%,\\
E_{Pn}&=\frac{100|-9.056-(-7.32)|}{|-7.32|}=\frac{100(1.736)}{7.32}\approx 23.72\%.
\end{aligned}
\]
\end{minipage}\par
\par\noindent\begin{minipage}{\linewidth}
\par\textit{Upper cutoff comparison.}
\[
E_A=|r_T-r_N|,\qquad E_{Pt}=\frac{100E_A}{|r_T|},\qquad E_{Pn}=\frac{100E_A}{|r_N|}.
\]
Here $r_T$ and $r_N$ are the unrounded contextual cutoffs being compared.
\[
\begin{aligned}
E_A&=|13.056-(11.32)|=1.736\approx 1.74\ \text{kg},\\
E_{Pt}&=\frac{100|13.056-(11.32)|}{|13.056|}=\frac{100(1.736)}{13.056}\approx 13.30\%,\\
E_{Pn}&=\frac{100|13.056-(11.32)|}{|11.32|}=\frac{100(1.736)}{11.32}\approx 15.34\%.
\end{aligned}
\]
\end{minipage}\par

\par\noindent\begin{minipage}{\linewidth}
\par\textbf{Step 7: Interpretation.}\nopagebreak\enspace The central $98\%$ limits are $(-9.06,13.06)$ kg under $t$ and $(-7.32,11.32)$ kg under normal, with approximately $1\%$ below each lower limit and $1\%$ above each upper limit. Lower limits mean mass below target; upper limits mean mass above target. Each normal endpoint is about $1.74$ kg closer to the mean deviation, so the normal interval is narrower.
\end{minipage}\par
\end{enumerate}
\endgroup
\end{document}
